% GENERATED FILE. Edit canonical lessons, then run npm run generate:books.
% canonical-source-hash: fnv1a64:1719211da989682c
% canonical-lessons: ES-C731-maduro, ES-C731-madurez, ES-C731-madurar, ES-C731-optimista, ES-C731-pesimista

\chapter{Crecer por dentro}
\label{ch:es-crecer-por-dentro}
\hlchaptermodality{731}

\begin{chapteropening}
\textbf{By the end of this chapter:} \emph{I can talk about maturity and whether I see life with hope or gloom.}

Complete the last lesson of chapter 731: I can talk about maturity and whether I see life with hope or gloom.
\end{chapteropening}

\section[maduro]{maduro — mature; ripe}

\label{lesson:ES-C731-maduro}



\begin{quote}
Before the new one: say the Spanish for beauty, then the Spanish for elegance.
\end{quote}

\subsection*{You'll want to know: maduro}
\textbf{maduro} — \textquotedblleft{}mature; ripe\textquotedblright{}.

\begin{quote}
Es muy maduro para su edad.
\end{quote}

\begin{quote}
\emph{He is very mature for his age.}
\end{quote}

With \emph{estar} it is ripe fruit: \emph{el plátano está maduro}.

\begin{cousinweb}
\emph{maduro} is inherited from Latin \emph{mātūrus} \textquotedblleft{}ripe, mature\textquotedblright{} (EtymDB); the Latin \emph{-t-} softened to \emph{-d-}.
\end{cousinweb}

\begin{grammarlens}[title={What you've built}]
One more word for telling, explaining and describing at B1.
\end{grammarlens}

\subsection*{Your turn}
\begin{itemize}
\raggedright
  \item \emph{Say it:} \emph{maduro}
  \item \emph{Say it:} \emph{maduro}, once more
  \item \emph{Say it:} a sentence of your own with \emph{maduro}
  \item \emph{Recall:} say the Spanish for beauty, then the Spanish for elegance, then say \emph{maduro} again
\end{itemize}

\begin{tcolorbox}[breakable,colback=teal!4,colframe=teal!35!black,title={Before you move on}]
What does maduro mean? (\textquotedblleft{}Mature; ripe\textquotedblright{}.) Say it once more.
\end{tcolorbox}
\section[la madurez]{la madurez — maturity}

\label{lesson:ES-C731-madurez}



\begin{quote}
Before the new one: say the Spanish for elegance, then the Spanish for mature.
\end{quote}

\subsection*{You'll want to know: la madurez}
\textbf{la madurez} — \textquotedblleft{}maturity\textquotedblright{}.

\begin{quote}
Actuó con madurez.
\end{quote}

\begin{quote}
\emph{She acted with maturity.}
\end{quote}

\begin{cousinweb}
\emph{madurez} is built on \emph{maduro} with the ending \emph{-ez}; \emph{maduro} is inherited from Latin \emph{mātūrus} (EtymDB).
\end{cousinweb}

\begin{grammarlens}[title={What you've built}]
One more word for telling, explaining and describing at B1.
\end{grammarlens}

\subsection*{Your turn}
\begin{itemize}
\raggedright
  \item \emph{Say it:} \emph{la madurez}
  \item \emph{Say it:} \emph{la madurez}, once more
  \item \emph{Say it:} a sentence of your own with \emph{la madurez}
  \item \emph{Recall:} say the Spanish for elegance, then the Spanish for mature, then say \emph{la madurez} again
\end{itemize}

\begin{tcolorbox}[breakable,colback=teal!4,colframe=teal!35!black,title={Before you move on}]
What does la madurez mean? (\textquotedblleft{}Maturity\textquotedblright{}.) Say it once more.
\end{tcolorbox}
\section[madurar]{madurar — to mature, to ripen}

\label{lesson:ES-C731-madurar}



\begin{quote}
Before the new one: say the Spanish for mature, then the Spanish for maturity.
\end{quote}

\subsection*{You'll want to know: madurar}
\textbf{madurar} — \textquotedblleft{}to mature, to ripen\textquotedblright{}.

\begin{quote}
Has madurado mucho este año.
\end{quote}

\begin{quote}
\emph{You have grown up a lot this year.}
\end{quote}

For people and fruit.

\begin{cousinweb}
\emph{madurar} comes from Latin \emph{mātūrāre} \textquotedblleft{}to ripen\textquotedblright{} (EtymDB), built on \emph{mātūrus}.
\end{cousinweb}

\begin{grammarlens}[title={What you've built}]
One more word for telling, explaining and describing at B1.
\end{grammarlens}

\subsection*{Your turn}
\begin{itemize}
\raggedright
  \item \emph{Say it:} \emph{madurar}
  \item \emph{Say it:} \emph{madurar}, once more
  \item \emph{Say it:} a sentence of your own with \emph{madurar}
  \item \emph{Recall:} say the Spanish for mature, then the Spanish for maturity, then say \emph{madurar} again
\end{itemize}

\begin{tcolorbox}[breakable,colback=teal!4,colframe=teal!35!black,title={Before you move on}]
What does madurar mean? (\textquotedblleft{}To mature, to ripen\textquotedblright{}.) Say it once more.
\end{tcolorbox}
\section[optimista]{optimista — optimistic; the optimist}

\label{lesson:ES-C731-optimista}



\begin{quote}
Before the new one: say the Spanish for maturity, then the Spanish for to mature.
\end{quote}

\subsection*{You'll want to know: optimista}
\textbf{optimista} — \textquotedblleft{}optimistic; the optimist\textquotedblright{}.

\begin{quote}
Soy optimista sobre el futuro.
\end{quote}

\begin{quote}
\emph{I am optimistic about the future.}
\end{quote}

Same form for men and women.

\begin{cousinweb}
\emph{optimista} is usually traced to French \emph{optimiste}, built on Latin \emph{optimus} \textquotedblleft{}best\textquotedblright{}.
\end{cousinweb}

\begin{grammarlens}[title={What you've built}]
One more word for telling, explaining and describing at B1.
\end{grammarlens}

\subsection*{Your turn}
\begin{itemize}
\raggedright
  \item \emph{Say it:} \emph{optimista}
  \item \emph{Say it:} \emph{optimista}, once more
  \item \emph{Say it:} a sentence of your own with \emph{optimista}
  \item \emph{Recall:} say the Spanish for maturity, then the Spanish for to mature, then say \emph{optimista} again
\end{itemize}

\begin{tcolorbox}[breakable,colback=teal!4,colframe=teal!35!black,title={Before you move on}]
What does optimista mean? (\textquotedblleft{}Optimistic; the optimist\textquotedblright{}.) Say it once more.
\end{tcolorbox}
\section[pesimista]{pesimista — pessimistic; the pessimist}

\label{lesson:ES-C731-pesimista}



\begin{quote}
Before the new one: say the Spanish for to mature, then the Spanish for optimistic.
\end{quote}

\subsection*{You'll want to know: pesimista}
\textbf{pesimista} — \textquotedblleft{}pessimistic; the pessimist\textquotedblright{}.

\begin{quote}
No seas tan pesimista.
\end{quote}

\begin{quote}
\emph{Do not be so pessimistic.}
\end{quote}

\begin{cousinweb}
\emph{pesimista} is usually traced to French \emph{pessimiste}, built on Latin \emph{pessimus} \textquotedblleft{}worst\textquotedblright{}.
\end{cousinweb}

\begin{grammarlens}[title={What you've built}]
One more word for telling, explaining and describing at B1.
\end{grammarlens}

\subsection*{Your turn}
\begin{itemize}
\raggedright
  \item \emph{Say it:} \emph{pesimista}
  \item \emph{Say it:} \emph{pesimista}, once more
  \item \emph{Say it:} a sentence of your own with \emph{pesimista}
  \item \emph{Recall:} say the Spanish for to mature, then the Spanish for optimistic, then say \emph{pesimista} again
\end{itemize}

\begin{tcolorbox}[breakable,colback=teal!4,colframe=teal!35!black,title={Before you move on}]
What does pesimista mean? (\textquotedblleft{}Pessimistic; the pessimist\textquotedblright{}.) Say it once more.
\end{tcolorbox}
